\section{Exploration des données}

\subsection{Chargement du jeu de données}

Le jeu de données  a été chargé au format .CSV à l’aide de la bibliothèque python \texttt{pandas} :

\begin{verbatim}
import pandas as pd
df = pd.read_csv("admission_data.csv")
\end{verbatim}

La commande \texttt{df.info()}permet de voir une partie des colonnes, de leurs types et du nombre de valeurs non nulles. On compte 9 colonnes et 550 lignes.

\subsection{Aperçu des données}

\begin{itemize}
    \item Le jeu de données contient des variables quantitatives comme par exemple (\textit{GRE Score}, \textit{TOEFL Score}, \textit{CGPA}, \textit{Chance of Admit}) et discrètes (\textit{University Rating}, \textit{SOP}, \textit{LOR}, \textit{Research}).
    \item Ici \textit{Chance of Admit} est la variable cible. Elle indique la probabilité d’admission.
\end{itemize}

\subsection{Traitement des valeurs manquantes}

Des valeurs manquantes ont été détectées dans deux colonnes :
\begin{itemize}
    \item \textbf{TOEFL Score}
    \item \textbf{SOP (Statement of Purpose)}
\end{itemize}

Ces deux colonnes présentaient chacune 37 valeurs manquantes.

La stratégie choisie a été de supprimer toutes les lignes contenant au moins une valeur manquante :

\begin{verbatim}
df_clean = df.dropna()
\end{verbatim}

Puis on a fait une vérification pour s'il n'y a plus de variable manquantes dans le dataset

\begin{verbatim}
df_clean.isnull().sum()
\end{verbatim}

Résultat : toutes les valeurs manquantes ont été supprimées. Le jeu de données final contient 513 lignes.

\subsection{Détection de doublons}

Pour vérifier la présence de doublons, j’ai utilisé la commande suivante :

\begin{verbatim}
df_clean.duplicated().sum()
\end{verbatim}

Résultat : aucune ligne dupliquée n’a été détectée.

\subsection{Étude des valeurs aberrantes}

Des diagrammes en boîte (\textit{boxplots}) ont été tracés pour les variables numériques afin de visualiser les potentiels outliers. Bien que certains points soient situés à l’extérieur du graphe de boites a moustache, aucun n’a été jugé significativement aberrant ou incohérent (par exemple, un CGPA supérieur à 10 ou un score GRE supérieur à 340). Aucun traitement particulier n’a été effectué à ce niveau.

\subsection{Étude des corrélations}

Une matrice de corrélation a été générée pour évaluer les relations entre les variables :

\begin{verbatim}
corr_matrix = df_clean.corr()
\end{verbatim}

Les coefficients de corrélation les plus élevés avec la variable cible \textit{Chance of Admit} sont :
\begin{itemize}
    \item \textbf{CGPA} : 0.87
    \item \textbf{GRE Score} : 0.82
    \item \textbf{TOEFL Score} : 0.79
    \item \textbf{Research} : 0.55
\end{itemize}

On voit que les performances académiques (CGPA, GRE, TOEFL) et l'expérience de recherche sont fortement corrélés aux chances d'admission.

\subsection{Conclusion de l'exploration}

L'analyse exploratoire a permis de :
\begin{itemize}
    \item Nettoyer le jeu de données (suppression des lignes incomplètes),
    \item Confirmer l'absence de doublons,
    \item Identifier les variables les plus corrélées à la cible,
    \item Préparer les données pour les étapes de visualisation et de modélisation.
\end{itemize}
